\documentclass[10pt,a4paper]{amsart}
\usepackage[margin=19mm]{geometry}
\usepackage{amsmath,amssymb,mathtools}
\usepackage[scr=rsfs,cal=euler]{mathalfa}
\usepackage{xfrac}
\usepackage[unicode,hidelinks,bookmarks=false]{hyperref}
\newcommand{\ie}{i.e.\ }
\newcommand{\cf}{cf.\ }
\newcommand{\resp}{respectively\ }
\newcommand{\Subsection}{\subsection}
\providecommand{\nobreakdash}{\nobreak\hbox{-}}
\setlength{\emergencystretch}{2em}
\multlinegap0pt
\usepackage{mathtools}
\usepackage{amssymb}
\usepackage[shortlabels]{enumitem}
\usepackage{tikz}
\usepackage{cancel}
\newtheorem{lemma}{Lemma}
\newcommand{\diver}{\operatorname{div}}
\newcommand{\tr}{\operatorname{tr}}
\title{Foliations by Critical Surfaces of the Hawking Energy in Asymptotically Flat Initial Data Sets}
\author{Alejandro Pe\~nuela Diaz}
\date{Source: arXiv:2508.13782v2 (CC BY 4.0)}
\begin{document}
\maketitle

\noindent\emph{Source and licence.} Foliations by Critical Surfaces of the Hawking Energy in Asymptotically Flat Initial Data Sets. Version arXiv:2508.13782v2 (CC BY 4.0), distributed by arXiv under the Creative Commons Attribution 4.0 International licence, http://creativecommons.org/licenses/by/4.0/, as recorded in the arXiv licence metadata for this exact version and shown on the abstract page https://arxiv.org/abs/2508.13782v2. Full licence notice: \texttt{licenses/arxiv-2508.13782-CC-BY-4.0.txt}.\par\smallskip

\noindent\textit{Editorial scope.} Counted unit: the article's lemma eulag1 (source0.tex lines 533--540). The article's complete original proof of it is delivered with it (source0.tex lines 541--601, one \texttt{proof} environment of 61 lines). No mathematics is written, repaired or re-derived: the statement and the proof are the article's own lines, in the article's own notation, and no retained line is substituted or re-flowed.  No further source-local context is needed: every cross-reference of every retained line resolves inside the two delivered spans.  Nothing was omitted from any retained span, so the package carries no omission record.  No retained line contains a citation, so the wrapper carries no bibliography and no bibliography entry.  No retained line contains a figure environment, an image-inclusion command or drawing code, so no image is delivered and the package declares no asset dependency.  Editorial formatting only, in addition: an AMS \texttt{amsart} preamble that restates the article's own declarations and macros used by the retained lines, namely the environments \texttt{lemma} on separate counters, together with the standard packages those lines need.  The retained environments print in this document as Lemma 1 in order of appearance, whereas the article prints the same items with the numbering of its own full document.  The complete native source of the article as distributed in the arXiv e-print for this exact version is bundled unchanged as \texttt{sources/183508/source0.tex}.
\begin{lemma}
There exist open neighborhoods $\mathcal{U}_1$ of $\delta \in  \mathcal{G}_1$, $\mathcal{U}_2$ of $0 \in  \mathcal{G}_2$,  $\mathcal{V}$ of $0 \in \Lambda_1(S_1(0))^\bot$, and $I$ of $0 \subset \mathbb{R}$ as
well as smooth maps $u : \mathcal{U}_1 \times \mathcal{U}_2 \longrightarrow \mathcal{V}$ and $\lambda : \mathcal{U}_1 \times \mathcal{U}_2 \longrightarrow I$ such that for every $(g,k)\in\mathcal U_1\times\mathcal U_2$ the surface $\Sigma_{0,1}(u(g,k))$ has area equal to $4 \pi$ and satisfies
\begin{equation}\label{eulag1}
 \lambda H  +\Delta^\Sigma H + H|\mathring{B}|^2+ H \mathrm{Ric}(\nu, \nu)+P( \nabla_\nu \tr k - \nabla_\nu k(\nu,\nu )) - 2\diver_\Sigma (Pk(\cdot, \nu)) +\frac{1}{2}H P^2 \in \Lambda_1(S_1(0)).
\end{equation}  
Here, all geometric quantities are computed with respect to the surface $\Sigma_{0,1}(u(g))$, the metric $g$, and the tensor $k$. Moreover, if $(g_0,k_0) \in \mathcal{U}_1 \times \mathcal{U}_2 $,  $u_0 \in \mathcal{V} $ and $\lambda_0 \in I$ are such that $\Sigma_{0,1}(u_0)$ satisfies (\ref{eulag1}) with respect to $g_0$ and $k_0$, and has area equal to $4 \pi$, then $u_0 = u(g_0, k_0)$ and $\lambda_0 = \lambda (g_0, k_0)$.
\end{lemma}

\begin{proof}
The argument follows the same general scheme as in the time-symmetric case (\(k=0\)),
but since it constitutes a central step in the analysis we present a complete version here for clarity.

   Let $\Tilde{\mathcal{V}} $ be a neighborhood of $0 \in \Lambda_1(S_1(0))^\bot $,  $\Tilde{\mathcal{U}}_1 $ a neighbourhood of the Euclidean metric $\delta \in  \mathcal{G}_1$, and $\Tilde{\mathcal{U}}_2 $ a neighborhood of $0 \in  \mathcal{G}_2$, such that the map 
   \begin{equation}
       T: \Tilde{\mathcal{V}} \times \mathbb{R} \times \Tilde{\mathcal{U}}_1 \times \Tilde{\mathcal{U}}_2   \longrightarrow \Lambda_1(S_1(0))^\bot \times  \mathbb{R}
   \end{equation}
   given by 
   \begin{equation}
       T(u,\lambda, g, k)= \left( P_{\Lambda_1(S_1(0))^\bot}(W_1(g)+H \lambda +W_2(g,k)), |\Sigma|   \right)
   \end{equation}
where $P_{\Lambda_1(S_1(0))^\bot}$ denotes the orthogonal projection to $ \Lambda_1(S_1(0))^\bot$, is well-defined and smooth.
   
The derivative in the direction of the first variable at point $(0,0,\delta, 0) $ is 
$$(dT)_{(0,0,\delta, 0)} ( (u,0, 0,0)) = (  -\Delta^{\mathbb{S}^2} (- \Delta^{\mathbb{S}^2} -2)u   ,\,  8 \pi P_{\Lambda_0(S_1(0))} u  ) $$
and the derivative in the direction of the second variable at point $(0,0,\delta, 0) $ is 
$$(dT)_{(0,0,\delta, 0)} ( (0,\lambda, 0,0)) = (2 \lambda, 0 ) $$
Then the differential $dT$ is 
\begin{equation}
dT(u,\lambda) = 
    \begin{pmatrix}
 -\Delta^{\mathbb{S}^2} (- \Delta^{\mathbb{S}^2} -2)  & 2  \\
 8 \pi P_{\Lambda_0(S_1(0))}  & 0 
\end{pmatrix}
\begin{pmatrix}
u \\
\lambda \\
\end{pmatrix}
\end{equation}
If we manage to see that $dT$ is an isomorphism, then by the implicit function theorem, we would have the result.  Now to see the $dT$ is injective, suppose that there are $u_1,u_2 \in  \Tilde{\mathcal{V}} \subset \Lambda_1(S_1(0))^\bot $ and $ \lambda_1, \lambda_2 \in \mathbb{R} $ such that 
\begin{equation}
    \begin{pmatrix}
 -\Delta^{\mathbb{S}^2} (- \Delta^{\mathbb{S}^2} -2)  & 2  \\
 8 \pi P_{\Lambda_0(S_1(0))}  & 0 
\end{pmatrix}
\begin{pmatrix}
u_1 \\
\lambda_1 \\
\end{pmatrix}
=
\begin{pmatrix}
 -\Delta^{\mathbb{S}^2} (- \Delta^{\mathbb{S}^2} -2)  & 2  \\
 8 \pi P_{\Lambda_0(S_1(0))}  & 0 
\end{pmatrix}
\begin{pmatrix}
u_2 \\
\lambda_2 \\
\end{pmatrix}
\end{equation}
this implies 
$$  8 \pi P_{\Lambda_0(S_1(0))}u_1=  8 \pi P_{\Lambda_0(S_1(0))}u_2, \quad    -\Delta^{\mathbb{S}^2} (- \Delta^{\mathbb{S}^2} -2) u_1 +2\lambda_1=  -\Delta^{\mathbb{S}^2} (- \Delta^{\mathbb{S}^2} -2) u_2 +2\lambda_2  $$

Combining the first equation with the second we obtain.
$$   -\Delta^{\mathbb{S}^2} (- \Delta^{\mathbb{S}^2} -2) (\Tilde{u}_1- \Tilde{u}_2) = 2(\lambda_2 -\lambda_1) $$
where we denote $\Tilde{u}:=P_{(\Lambda_0(S_1(0)) \times \Lambda_1(S_1(0)))^\bot}(u) $ the orthogonal projection of $u$ to $(\Lambda_0(S_1(0)) \times \Lambda_1(S_1(0)))^\bot$, we also denote $\Lambda_i(S_1(0))$  by $\Lambda_i$ .

 Since $-\Delta^{\mathbb{S}^2} (- \Delta^{\mathbb{S}^2} -2):(\Lambda_0 \times \Lambda_1)^\bot \to (\Lambda_0 \times \Lambda_1)^\bot $ is an isomorphism and $\Tilde{u}_1- \Tilde{u}_2 \in (\Lambda_0 \times \Lambda_1)^\bot $, then $2(\lambda_2 -\lambda_1) \in (\Lambda_0 \times \Lambda_1)^\bot$, and with this we have  that $\lambda_1= \lambda_2$ and $u_1 =u_2 $. 

 Now to show that  $dT$ is surjective  consider a given  $u_1 \in  \Tilde{\mathcal{V}} \subset \Lambda_1^\bot $ and $ \lambda_1 \in \mathbb{R} $.  By choosing  $u_2 \in  \Tilde{\mathcal{V}} \subset \Lambda_1^\bot $ and $ \lambda_2 \in \mathbb{R} $ such that  $\lambda_2 = \frac{1}{2} P_{\Lambda_0}(u_1) $,  $P_{(\Lambda_0 \times \Lambda_1)^\bot}(u_2)=(-\Delta^{\mathbb{S}^2} (- \Delta^{\mathbb{S}^2} -2))^{-1} P_{(\Lambda_0 \times \Lambda_1)^\bot}(u_1) $ and  $P_{\Lambda_0}(u_1)=P_{\Lambda_0}(u_1)$, it holds $dT(u_2,\lambda_2)= (u_1,\lambda_1)^T$. Then $dT$ is an isomorphism and by the implicit function theorem we have the result.
\end{proof}

\end{document}
